% Tipado autónomo de la carta dodecafásica terminal.
\subsection{Carta dodecafásica firmada del estado TPK pleno}
\label{subsec:registro-evaluacion-incidencial}

La carta dodecafásica y una eventual factorización por la proyección cruda de
108 pasos (CT108) son objetos de tipos distintos. Sea
\(X\in\mathsf{TPKFullState}\) un estado enriquecido completo y sea
\(\rho_{108}(X)\in\mathsf{CT108RawProjection}\) el objeto obtenido al olvidar
orientación, hoja, acarreo, firma de frontera, memoria y libro dodecafásico.
En las coordenadas canónicas del estado pleno escribimos
\[
 X=(X_{\mathrm{APP}},X_{\mathrm{TRIT}},X_{\mathrm{ruta}},X_{\mathrm{mem}},
       C_{\mathrm{term}},U_{12}^{\mathrm{sgn}},
       X_{\mathrm{frontera}},\ldots).
\]
El lector activo es, por definición, la proyección de esa coordenada ya
generada y conservada:
\[
 \mathcal S_{12}:=\operatorname{pr}_{U_{12}^{\mathrm{sgn}}}:
 \mathsf{TPKFullState}\longrightarrow\mathsf{SignedDodecaphaseState}.
\]
En la subred compatible del estado pleno, esta proyección coincide con el
lector compuesto
\begin{equation}
 \mathcal S_{12}\equiv R_{\mathrm{sgn}}
 :=\Pi_H\circ\mathcal E_{108}^{90,120}\circ\mathcal R_{12}.
 \label{eq:registro-identidad-s12-rsgn}
\end{equation}
Por tanto, \(\mathcal S_{12}\) lee una coordenada conservada del mismo estado;
no recibe una constante
convencional ni una tabla metrológica. No se sustituye su dominio por
\(\mathsf{CT108RawProjection}\), pues esa sustitución exigiría además una
factorización \(\mathcal S_{12}=\Xi\circ\rho_{108}\). Tal factorización no se
presupone ni se necesita aquí; requeriría demostrar que la lectura firmada es
constante sobre las fibras de \(\rho_{108}\). El teorema interno opera
directamente sobre el estado TPK pleno.

Para fijar el contenido de una realización incidencial posible, sea
\(\mathscr L\) un libro de visitas y trazas orientadas perteneciente a
\(X\). Cada visita conserva la ruta, la posición APP, la hoja, el
instante, la multiplicidad y la incidencia de frontera. La evaluación
aritmética produce el entero y su descomposición
\[
 n_v=r_v+9q_v,\qquad 1\leq r_v\leq9.
\]
El entero \(n_v\) se calcula a partir de la hoja de suma o producto; su
residuo no se recibe como un valor independiente. La distinción entre
incidencia de corona e incidencia interior permanece explícita para las
visitas de residuo \(9\).

En la ventana temporal
\(E_m=\{9(m-1)+1,\ldots,9m\}\), se distinguen los dominios
\[
 \mathscr V_m=\{v:t(v)\in E_m\},\qquad
 \mathscr T_m=\{\theta:\operatorname{origen}(\theta)\in E_m\}.
\]
La presentación por recuentos considera
\begin{align}
 A_m&=\sum_{v\in\mathscr V_m}w(v)
             \mathbf1_{\{1,3,5,7\}}(r(v)),\\
 C_m&=-\sum_{\theta\in\mathscr T_m}u(\theta)\sigma(\theta)
      =\sum_{j\geq0}\bigl(\nu^-_{120,m}(j)-\nu^+_{120,m}(j)\bigr),\\
 V_m&=\sum_{v\in\mathscr V_m}w(v)
              \mathbf1_{\{9\}}(r(v))\epsilon_\partial(v).
\end{align}
Aquí \(w(v)\) es la multiplicidad incidencial y
\(\epsilon_\partial(v)=1\) en la corona, \(-1\) en el interior.
Las trazas contadas por \(\nu^\pm\) deben enumerarse mediante la regla
correspondiente del libro. La suma entera requiere soporte finito o una
construcción alternativa que determine su significado. El número de
incidencias cronológicas y la longitud reducida \(120(1+3j)\) tienen
tipos distintos; su relación exige la unidad de longitud del lector.

Se define \([x]_{10}\in\{0,\ldots,9\}\) como el representante
euclídeo y se conservan los cocientes
\[
 x=[x]_{10}+10\lfloor x/10\rfloor.
\]
El bloque incidencial es
\begin{equation}
 z_m=100[A_m]_{10}+10[C_m]_{10}+[V_m]_{10}.
 \label{eq:registro-bloque-incidencial}
\end{equation}
Por composición se obtiene la aplicación
\begin{equation}
 \mathcal E_{\rm cnt}(\mathscr L)=
 \left((z_m-z_{m+3})_{m=1}^{12},
       (z_m-z_{m+4})_{m=1}^{12},\sum_{m=1}^{12}z_m\right).
\label{eq:registro-composicion-incidencial}
\end{equation}
Los índices se leen cíclicamente módulo \(12\). En este apartado,
\((D_jz)_m=z_m-z_{m+j}\), para \(j=3,4\), y
\(Q(z)=\sum_{m=1}^{12}z_m\); por tanto, la composición anterior es
\((D_3z,D_4z,Q(z))\). Esta convención se conserva en
\S\ref{subsec:k-transversal}, donde se demuestra la reconstrucción
integral completa.
Esta fórmula produce sus coordenadas a partir de los recuentos. No utiliza
\(K\), \(U\), \(\alpha\) ni las coordenadas transversales esperadas
como argumentos.

\subsubsection{Publicación terminal y cierre exacto de sus cartas}

El estado terminal canónico es
\[
 X_{\mathrm{term}}
 =\bigl(\mathcal U_{108}\circ\cdots\circ\mathcal U_1\bigr)
   (X_{\mathrm{seed}}),
 \qquad
 \mathcal U_t=\mathsf{Upd}_t\circ\mathsf{Tra}_t\circ\mathsf{Sel}_t,
\]
donde cada \(\mathcal U_t\) actúa sobre el estado TPK pleno y no sobre su
proyección cruda. Entre las coordenadas producidas y preservadas antes de la
publicación firmada se encuentra
\[
 C_{\mathrm{term}}=\operatorname{pr}_{C}(X_{\mathrm{term}})
 =(b^{90},b^{120},Q_{\mathrm{TPK}}),
\]
con
\begin{align*}
 b^{90}={}&(-495,-116,-684,108,601,-90,
              -173,-88,313,560,-397,461),\\
 b^{120}={}&(-425,-281,-481,671,-255,30,
               475,-543,680,251,6,-128),\\
 Q_{\mathrm{TPK}}={}&6263.
\end{align*}
No se introduce aquí una tabla posterior: \(C_{\mathrm{term}}\) es la
coordenada transversal que \(\mathcal E_{108}^{90,120}\) produce dentro de
la evolución del estado pleno. Sobre la subred compatible de estados TPK la
carta satisface la identidad de operadores
\begin{equation}
 \mathcal S_{12}=\operatorname{pr}_{U}
 =\Pi_H\circ\operatorname{pr}_{C}.
 \label{eq:registro-identidad-s12-pih}
\end{equation}
La identidad se formula sobre el estado TPK pleno y no presupone el descenso
de \(\mathcal S_{12}\) por \(\rho_{108}\).

La evaluación posterior desde los veinticinco campos publicados queda
expuesta sin salto. Para \(1\le r\le3\),
sean
\[
 B_r=\sum_{j=0}^{3}b^{120}_{r+3j},\qquad
 d_{r,j}=b^{90}_{r+3j},
\]
y
\[
 s_1=\frac{Q_{\mathrm{TPK}}+2B_1+B_2}{3},\qquad
 s_2=s_1-B_1,\qquad s_3=s_2-B_2.
\]
La sustitución de los veinticinco campos anteriores da, enteramente,
\[
 (B_1,B_2,B_3)=(972,-1073,101),\qquad
 (s_1,s_2,s_3)=(2378,1406,2479).
\]
Aplicando
\[
 U^{(r)}=(s_r,d_{r,1}-d_{r,3},d_{r,0}+d_{r,2},d_{r,0}-d_{r,2})
\]
se obtienen, componente a componente,
\begin{align*}
 U^{(1)}&=(2378,-452,-668,-322),\\
 U^{(2)}&=(1406,998,-204,-28),\\
 U^{(3)}&=(2479,-551,-371,-997).
\end{align*}
La intercalación por columnas evalúa por tanto la carta firmada:
\begin{equation}
 U_{\mathrm{term}}:=\mathcal S_{12}(X_{\mathrm{term}})
 =(2378,1406,2479,-452,998,-551,-668,-204,-371,-322,-28,-997).
 \label{eq:registro-uterm-desde-estado-pleno}
\end{equation}

\begin{theorem}[Cierre exacto posterior entre \texorpdfstring{\(U\)}{U}, \texorpdfstring{\(K\)}{K} y los canales]
\label{thm:registro-evaluacion-terminal-recuperada}
Una vez publicada por proyección la carta firmada
\eqref{eq:registro-uterm-desde-estado-pleno}, la transformada orbital de
Hadamard publica, en una etapa posterior,
\[
 K=(234,543,140,729,659,824,621,058,914,794,146,601),
\]
y el extractor transversal \((D_3,D_4,Q)\) produce
\begin{align}
 b^{90}={}&(-495,-116,-684,108,601,-90,
             -173,-88,313,560,-397,461),\notag\\
 b^{120}={}&(-425,-281,-481,671,-255,30,
              475,-543,680,251,6,-128),\notag\\
 Q_{\mathrm{TPK}}={}&6263.
 \label{eq:registro-canales-terminales-producidos}
\end{align}
La composición es una biyección exacta entre la subred compatible de estados
firmados, el sello \(K\) y su coordenada transversal
\(C_{\mathrm{term}}=(b^{90},b^{120},Q_{\mathrm{TPK}})\).
En particular,
\begin{equation}
 C_{\mathrm{term}}=(D_3K,D_4K,Q(K)).
 \label{eq:registro-cterm-desde-xterm}
\end{equation}
\end{theorem}

\begin{proof}
Agrupando \(U_{\mathrm{term}}\) en sus tres órbitas de paso tres se obtiene
\[
 \begin{aligned}
  u^{(0)}&=(2378,-452,-668,-322),&
  \tfrac14H_4u^{(0)}&=(234,729,621,794),\\
  u^{(1)}&=(1406,998,-204,-28),&
  \tfrac14H_4u^{(1)}&=(543,659,58,146),\\
  u^{(2)}&=(2479,-551,-371,-997),&
  \tfrac14H_4u^{(2)}&=(140,824,914,601).
 \end{aligned}
\]
Las divisiones son exactas en \(\mathbb Z\).
La intercalación por columnas publica \(K\). Finalmente, las doce restas
cíclicas de pasos tres y cuatro y la suma total dan, componente a componente,
\eqref{eq:registro-canales-terminales-producidos}. La proposición
\ref{prop:k-transversal} demuestra la inversa integral, y el
teorema~\ref{thm:k-naturalidad-lector-firmado} demuestra la naturalidad bajo
refinamiento. Ninguna operación recibe \(\alpha\), CODATA, una tabla
metrológica ni una constante convencional.
\end{proof}

\begin{remark}[Forma operatoria completa y control ejecutable focal]
La flecha terminal no se postula a partir de una tabla: es la composición
tipada
\[
 X_{\mathrm{term}}
 \xrightarrow{\mathcal R_{12}}\mathscr L(X_{\mathrm{term}})
 \xrightarrow{\mathcal E_{108}^{90,120}}C_{\mathrm{term}}
 \xrightarrow{\Pi_H}U_{\mathrm{term}}
 \xrightarrow{\mathcal H_{12}^{-1}}K .
\]
Las definiciones anteriores recorren las doce ventanas que particionan los
ciento ocho pasos y expresan cada coordenada de canal como una suma finita de
sus visitas y trazas, con orientación, hoja, acarreo, frontera y memoria. Ésta
es la materialización matemática de la producción de
\(C_{\mathrm{term}}\); una expansión tabular fila por fila sería otra
serialización de la misma suma, no una premisa adicional del teorema. El
ejecutable focal reevalúa exactamente la carta finita
\(C_{\mathrm{term}}\leftrightarrow U_{\mathrm{term}}\leftrightarrow K\)
como control reproducible suplementario. La demostración de la hipótesis del
continuo permanece formulada sobre APP--TRIT--TPK, el estado enriquecido
completo, sus publicaciones correlacionadas y la estructura discreta conjunta.
\end{remark}

Ni \(\mathcal S_{12}\), ni \(U\), ni \(K\), ni \(\alpha\) son entradas
externas ni selectores convencionales de la construcción discreta del
continuo. Son coordenadas y publicaciones tipadas del mismo estado
APP--TRIT--TPK. La coordenada \(\alpha_{\mathrm{HMT}}\) es una salida derivada
y correlacionada de esa generación única; su orden de publicación no rompe la
continuidad del objeto ni debilita su derivación.

La fórmula incidencial \eqref{eq:registro-composicion-incidencial} describe
una realización adicional cuando se despliega un libro completo
\(\mathscr L(X_{\mathrm{term}})\). El teorema anterior no afirma que la
coordenada firmada pueda reconstruirse después de olvidar ese libro mediante
\(\rho_{108}\); tampoco necesita tal afirmación. Esta separación evita
confundir el cierre interno del estado enriquecido con un problema inverso
planteado sobre una proyección que ha eliminado precisamente la información
de orientación y memoria.

La aplicación \(\Pi_H\), ejecutada dentro de la carta, produce a partir de
\(C_{\mathrm{term}}\) la coordenada firmada y sus tres órbitas de paso tres:
\begin{align*}
 U_{\mathrm{term}}^{(1)}&=\Pi_H^{(1)}(C_{\mathrm{term}})
                         =(2378,-452,-668,-322),\\
 U_{\mathrm{term}}^{(2)}&=\Pi_H^{(2)}(C_{\mathrm{term}})
                         =(1406,998,-204,-28),\\
 U_{\mathrm{term}}^{(3)}&=\Pi_H^{(3)}(C_{\mathrm{term}})
                         =(2479,-551,-371,-997).
\end{align*}
Sea \(\mathcal P_{\mathrm{orb}}:\mathbb Z^{12}\to(\mathbb Z^4)^3\) el
reagrupamiento orbital
\[
 \mathcal P_{\mathrm{orb}}(U)
 =\bigl((U_1,U_4,U_7,U_{10}),
        (U_2,U_5,U_8,U_{11}),
        (U_3,U_6,U_9,U_{12})\bigr),
\]
y sea \(\operatorname{Int}_{3,4}=\mathcal P_{\mathrm{orb}}^{-1}\) la
intercalación por columnas de tres cuaternas. Definimos además
\[
 \operatorname{Dig}_{10}^{(3)}(z)
 =\left(\left\lfloor\frac z{100}\right\rfloor,
         \left\lfloor\frac z{10}\right\rfloor\bmod10,
         z\bmod10\right),\qquad0\le z\le999.
\]
La carta modular terminal es la aplicación efectiva sobre el estado compatible
de canales
\begin{equation}
 \mathcal G_{\mathrm{term}}(C)
 :=\left[(\operatorname{Dig}_{10}^{(3)})^{12}
 \circ\operatorname{Int}_{3,4}
 \circ\left(\tfrac14H_4\oplus\tfrac14H_4
                    \oplus\tfrac14H_4\right)
  \circ\mathcal P_{\mathrm{orb}}\right]\!\left(\Pi_H(C)\right).
 \label{eq:registro-generador-modular-terminal}
\end{equation}
Su dominio es el subconjunto de la subred de integralidad de coordenadas
compatibles de \(\mathbb Z^{25}\)
en el que \(\Pi_H\) es integral, las tres imágenes orbitales por \(H_4\) son
divisibles por cuatro y las doce coordenadas resultantes pertenecen a
\(\{0,\ldots,999\}\). No recibe \(U_{\mathrm{term}}\), \(K\), \(\alpha\),
una tabla decimal ni un valor metrológico; la intercalación, el origen de fase
y la orientación forman parte del tipo de la coordenada de canales.

La coordenada de comprobación se serializa en
\path{metadata/semilla-registro-terminal.json}. Ese archivo contiene
\(C_{\mathrm{term}}\), el origen de fase y la orientación; no contiene
\(U_{\mathrm{term}}\), \(K\) ni las doce filas residuales. El programa focal
aplica \(\mathcal G_{\mathrm{term}}\), publica \(U_{\mathrm{term}}\), \(K\) y
los residuos, y retorna por \((D_3,D_4,Q)\) al mismo
\(C_{\mathrm{term}}\). Comenzar allí permite una comprobación focal exacta
de la carta finita posterior. No redefine el dominio de la carta firmada
\eqref{eq:registro-uterm-desde-estado-pleno}, no reemplaza la prueba precedente y no
se atribuye al auxiliar la producción de la coordenada que recibe.

\begin{proposition}[Cálculo modular de las doce filas terminales]
\label{prop:registro-generacion-modular-terminal}
La evaluación de \(\mathcal G_{\mathrm{term}}\) sobre la coordenada producida
\(C_{\mathrm{term}}\) obtiene primero \(U_{\mathrm{term}}\) y después las
cuaternas de fase
\[
 (234,729,621,794),\qquad(543,659,058,146),\qquad
 (140,824,914,601),
\]
y, tras la intercalación, produce exactamente las doce filas de la tabla
siguiente. Es la comprobación algebraica posterior al
teorema~\ref{thm:registro-evaluacion-terminal-recuperada}; no sustituye la
publicación de la carta firmada del estado TPK pleno.
\end{proposition}

\begin{proof}
Las fórmulas de \(\Pi_H\) dan los tres bloques firmados escritos arriba. El
producto por \(H_4/4\) da después, sin redondeo,
\begin{align*}
 \tfrac14H_4(2378,-452,-668,-322)^T&=(234,729,621,794)^T,\\
 \tfrac14H_4(1406,998,-204,-28)^T&=(543,659,58,146)^T,\\
 \tfrac14H_4(2479,-551,-371,-997)^T&=(140,824,914,601)^T.
\end{align*}
Las doce divisiones son exactas en \(\mathbb Z\). Intercalar por columnas da
\[
 \begin{split}
 (&234,543,140,729,659,824,\\
  &621,058,914,794,146,601),
 \end{split}
\]
y \(\operatorname{Dig}_{10}^{(3)}\) proporciona de forma única sus tres
residuos decimales. El verificador auxiliar incluido en el paquete ejecuta
estas mismas operaciones; sus controles negativos alteran el estado de canales
o su origen de fase.
\end{proof}

Denotemos por
\[
 \mathscr F_{\mathrm{term}}^{(10)}
 :=\bigl((a_m,c_m,v_m)\bigr)_{m=1}^{12}
\]
la carta de dígitos publicada por \(\operatorname{Dig}_{10}^{(3)}(K_m)\).
Cuando se dispone de una realización incidencial completa, se cumple
\((a_m,c_m,v_m)=([A_m]_{10},[C_m]_{10},[V_m]_{10})\); esa igualdad no se
postula para reconstruir el libro después de haberlo olvidado. Cada fila
conserva el índice de ventana y, por ello, no es una bolsa de cifras
intercambiables.

\begin{center}
\begin{tabular}{c@{\qquad}ccc@{\qquad}c}
\toprule
\(m\)&\(a_m\)&\(c_m\)&\(v_m\)&\(K_m\)\\
\midrule
1&2&3&4&234\\
2&5&4&3&543\\
3&1&4&0&140\\
4&7&2&9&729\\
5&6&5&9&659\\
6&8&2&4&824\\
7&6&2&1&621\\
8&0&5&8&058\\
9&9&1&4&914\\
10&7&9&4&794\\
11&1&4&6&146\\
12&6&0&1&601\\
\bottomrule
\end{tabular}
\end{center}
En consecuencia,
\begin{equation}
 z^{\mathrm{term}}
 =(234,543,140,729,659,824,621,058,914,794,146,601).
 \label{eq:registro-terminal-residuos}
\end{equation}

\begin{proposition}[Retorno terminal de la familia calculada a los canales]
\label{prop:registro-terminal-coordenadas}
La recomposición decimal de \(\mathscr F_{\mathrm{term}}^{(10)}\), seguida
del extractor transversal, produce, sin consultar \(\alpha\) ni una tabla
metrológica,
\begin{align*}
 D_3z^{\mathrm{term}}={}&(-495,-116,-684,108,601,-90,
 &\hspace{18mm}-173,-88,313,560,-397,461),\\
 D_4z^{\mathrm{term}}={}&(-425,-281,-481,671,-255,30,
 &\hspace{18mm}475,-543,680,251,6,-128),\\
 Q(z^{\mathrm{term}})={}&6263.
\end{align*}
Estas son exactamente las veinticinco coordenadas
\((b^{90},b^{120},Q_{\mathrm{TPK}})\) utilizadas por el lector armónico.
\end{proposition}

\begin{proof}
La formación decimal de cada fila da
\eqref{eq:registro-terminal-residuos}. Restar cíclicamente la entrada
\(m+3\) y la entrada \(m+4\) produce, componente a componente, los dos
vectores escritos; la suma de las doce entradas es \(6263\). El cálculo se
repite en el verificador autónomo del paquete. La familia terminal es una
publicación calculada del estado de canales. La evaluación transversal retorna
exactamente a la semilla \(C_{\mathrm{term}}\); la lectura armónica recupera
después \(U_{\mathrm{term}}\) y la inversa de Hadamard el mismo \(K\). Esto
prueba la concordancia de ambas coordenatizaciones sin usar alfa como selector.
\end{proof}

La carta \(\mathscr F_{\mathrm{term}}^{(10)}\) es suficiente y necesaria
para esta evaluación transversal. La tabla es un certificado decimal de la
salida publicada, no una reedición visita por visita del libro orientado. El
estado enriquecido conserva además cocientes euclídeos, visitas, trazas y
procedencia; esos campos no se declaran recuperables a partir de los treinta y
seis residuos, ni son necesarios para el teorema cardinal de este artículo.

\begin{proposition}[Información aritmética suficiente y necesaria]
\label{prop:registro-36-residuos}
Sean \(N,N'\in\mathbb Z^{12\times3}\), ordenados por los recuentos
\((A,C,V)\). Para la aplicación definida por
\eqref{eq:registro-bloque-incidencial}--\eqref{eq:registro-composicion-incidencial},
\[
 \mathcal E_{\rm cnt}(N)=\mathcal E_{\rm cnt}(N')
 \quad\Longleftrightarrow\quad
 N-N'\in10\mathbb Z^{36}.
\]
En particular, los \(36\) residuos sectoriales determinan exactamente
la salida de \(25\) coordenadas. La conservación de los cocientes y
de la procedencia constituye información adicional del registro.
\end{proposition}

\begin{proof}
La igualdad módulo \(10\) produce los mismos bloques y, por tanto,
las mismas coordenadas transversales. Recíprocamente, si las coordenadas
transversales coinciden, \(D_3(z-z')=D_4(z-z')=0\). Los pasos
\(3\) y \(4\) generan el ciclo de longitud \(12\), de modo que
\(z-z'\) es constante. La igualdad de cargas totales anula esa
constante. Finalmente, la representación decimal de un entero entre
\(0\) y \(999\) tiene tres dígitos únicos. Se recuperan así los
tres residuos de cada ventana. La afirmación es una identidad algebraica
para todo \(N,N'\), no una inferencia obtenida de pruebas finitas.
\end{proof}

\subsection{Conjugación de los recuentos y términos de frontera}
\label{subsec:registro-conjugacion-incidencial}

Sea \(\rho(m)=13-m\). Considérese una conjugación de incidencias
que invierte el orden de las ventanas, conserva la clase aritmética y la
clase de frontera de las visitas e intercambia los canales de las trazas.
Entonces
\[
 (A_m^\dagger,C_m^\dagger,V_m^\dagger)
 =(A_{\rho(m)},-C_{\rho(m)},V_{\rho(m)}).
\]
Si \(c_m=[C_m]_{10}\), la reducción decimal produce la identidad
\begin{equation}
 z_m^\dagger=z_{\rho(m)}
       +100\mathbf1_{\{c_{\rho(m)}\ne0\}}-20c_{\rho(m)}.
 \label{eq:registro-conjugacion-decimal}
\end{equation}
En efecto, \([-C]_{10}=0\) cuando \([C]_{10}=0\), y
\([-C]_{10}=10-[C]_{10}\) en otro caso. La fórmula conserva
exactamente esa diferencia. La negación del canal no equivale a negar
el bloque decimal completo.

La aplicación de esta identidad a una involución concreta del TPK exige
verificar su acción sobre las incidencias. Para un lector de ocupaciones
evaluado antes del avance, la inversión de una ruta
\(\gamma=(x_0,\ldots,x_n)\) satisface
\[
 \sum_{k=0}^{n-1}f(x_{n-k})
 =\sum_{k=0}^{n-1}f(x_k)+f(x_n)-f(x_0).
\]
Los términos de los extremos se incorporan antes de reducir módulo
\(10\). En rutas cerradas se anulan; en ventanas abiertas pueden
alterar sus residuos. Esta corrección, ya presente en el desarrollo
bidireccional del transporte, especifica qué debe conservar el empalme
entre la ruta y los recuentos.

\begin{proposition}[Necesidad de la información no periódica del registro]
Supóngase que el observable de una familia fija de rutas satisface
\(o(t+54)=o(t)\), que la multiplicidad se conserva bajo ese retorno
y que el mismo lector local actúa en las ventanas de \(9\) instantes.
Entonces sus recuentos y bloques satisfacen
\[
 (A,C,V)_{m+6}=(A,C,V)_m,\qquad z_{m+6}=z_m.
\]
\end{proposition}

\begin{proof}
La traslación \(t\mapsto t+54\) lleva \(E_m\) biyectivamente
a \(E_{m+6}\), conserva cada incidencia contada y su multiplicidad.
La igualdad se transmite a la suma y a la reducción decimal.
\end{proof}

Por tanto, una presentación con \(12\) bloques distintos requiere que
el lector conserve alguna información que no factorice por ese observable
de período \(54\): selección de incidencias, orientación efectiva,
extremos, multiplicidad o memoria. La proposición delimita una reducción
inadmisible del estado; no afirma que la dinámica completa tenga
período \(54\) ni determina por sí misma cuál es su lector terminal.
